\section{Formalization and Proof Ledger}

\subsection{Standalone hierarchy}

The machine-checked development separates generic coding mathematics from
application bridges:
\begin{leanledger}
\begin{description}[style=nextline]
  \item[\lean{Proximity/Core/BlockCode}]
    Relative distance, balls, and universal list-bound predicates.
  \item[\lean{Proximity/Core/ReedSolomon}]
    Evaluation-domain definitions, RS codes, and interleaving.
  \item[\lean{Proximity/Core/PlanePair}]
    Plane--pair--line equivalences.
  \item[\lean{Proximity/Core/FiniteThreshold}]
    Safe finite grids, greatest-safe and first-unsafe values, and adjacent
    safe/unsafe threshold certificates.
  \item[\lean{Proximity/Core/Distance} and \lean{Proximity/Core/RadiusShell}]
    Attained minimum distances, closed and strict nearness, and the exact
    half-open real intervals between consecutive Hamming-count shells.
  \item[\lean{Proximity/Core/PolynomialAgreement}]
    Polynomial multiplier certificates and the whole-pair agreement bound.
  \item[\lean{Proximity/Core/PolynomialAgreement/Boundary}]
    Normalized coefficient-system equivalence, kernel translates of fixed
    solution fibers, and upward closure of feasible response degrees.
  \item[\lean{Proximity/Core/PolynomialAgreement/Completeness} and
        \lean{Proximity/Core/PolynomialAgreement/Sharpness}]
    Interpolation on the disagreement set, exactness through half the
    minimum distance, and the sharp family on nonzero evaluation domains.
  \item[\lean{Proximity/Capacity/Parameters}]
    Explicit \(\epsstar\) and field-size inequalities.
  \item[\lean{Proximity/Capacity/Realization} and \lean{Proximity/Capacity/ListPhase}]
    Actual Galois fields, exact-order subgroup domains for every family
    member, and instantiated list/MCA grid and real-radius conclusions.
  \item[\lean{Proximity/ListDecoding/Capacity}]
    Agreement-set injection, the matching received-word construction, the
    universal endpoint equality, the fractional-target equivalence, and the
    next-shell coefficient-vector injection.
  \item[\lean{Proximity/ListDecoding/PhaseSplit}]
    All three fractional-target regimes, the exact endpoint-safe real
    interval, its excluded supremum, and the largest safe grid count.
  \item[\lean{Proximity/Capacity/UniqueRadius} and \lean{Proximity/Capacity/BoundaryList}]
    Exact uniqueness radius at every positive width, the one-codeword
    fractional budget criterion and real safe interval, and the boundary
    prime's count and real list thresholds.
  \item[\lean{Proximity/Capacity/AgreementPacking}]
    Finite-radius packing of all agreement subsets and the adjacent list
    threshold for budgets immediately below capacity.
  \item[\lean{Proximity/Capacity/JoThresholds}]
    The four length-\(64\) prime fields and smooth evaluation domains
    from Jo, with exact list thresholds at every positive width.
  \item[\lean{Proximity/Capacity/WidthReduction}]
    Scalar and interleaved list bounds agree when the number of candidate
    pairs in a violating list is smaller than the field, including the
    \(2^{-128}\) target for \(q<2^{256}\).
  \item[\lean{Proximity/Capacity/SubfieldListDescent}]
    Coefficient expansion is a count-ball isometry between RS codes
    over an extension and their base-field interleavings, with equal
    list cardinalities.
  \item[\lean{Proximity/Capacity/SubfieldThresholdArithmetic}]
    The length-\(128\) subgroup in \(\mathbb F_{641^{26}}\) and all integer
    comparisons for the paired threshold certificate in
    \cref{prop:subfield-crossings}, including the list safety theorem at
    count \(62\).
  \item[\lean{Proximity/MCA/Definitions}]
    Arbitrary-block-code same-set MCA definitions at count and relative
    radii.
  \item[\lean{Proximity/MCA/Threshold}]
    Exact worst bad-count spectrum and the equivalence between operational
    real error at \(r/n\) and the integer \(2^{-\lambda}\) budget.
  \item[\lean{Proximity/MCA/SpectrumCertificate}]
    Supplied-code grid-threshold pinning from a uniform upper bound at one
    radius and a concrete over-budget affine line at the next radius.
  \item[\lean{Proximity/MCA/SmallRadius}]
    Universal RS unit-spike and two-spike certificates, including
    \(B_{\mathcal C}(0)=1\) for every proper RS code.
  \item[\lean{Proximity/MCA/BoundaryField}]
    Primality, order-\(4096\) smooth domain, and the exact 128-bit radius
    transition for \(q=2^{128}+192513\).
  \item[\lean{Proximity/MCA/RealThreshold}]
    Count/real same-set event correspondence and the boundary-field safe
    interval \([0,1/4096)\), including supremum and nonattainment.
  \item[\lean{Proximity/MCA/Pair}]
    RS compressed interpolation certificates,
    bad-challenge injection, the uniform cardinality bound, and the
    explicit-family error corollary.
  \item[\lean{Proximity/Application/Calibration}]
    Generic source-validity model, pointwise distance non-collapse, and the
    uniform source-far to target-far implication.
  \item[\lean{Proximity/Application/DistanceCalibration}]
    Identification with attained scalar/pair code distance, interpolation
    covering bounds, and good-challenge farness with its exact boundary.
  \item[\lean{Proximity/Core/CanonicalLift} and \lean{Proximity/Core/Systematic}]
    Canonical real-lift spanning, actual systematic interpolation and
    integer/real syndrome compatibility. These are proved algebraic facts.
  \item[\lean{Proximity/Core/InterleavedMetric}]
    Native scalar-coordinate and whole-pair Hamming comparisons, retaining
    their different normalizations.
  \item[\lean{Proximity/Core/MetricSharpness}]
    The degree-\(k\) first column with zero second column attains factor two
    between pair and scalar-coordinate distance for proper RS codes.
  \item[\lean{Proximity/Core/DeepHole} and \lean{Proximity/Core/FixedChallenge}]
    Exact scalar/pair/flattened capacity distances for monomial received words,
    and capacity-far counterexamples to arbitrary fixed-challenge compression.
  \item[\lean{Proximity/MCA/ZeroDirection} and \lean{Proximity/MCA/Diagonal}]
    Explicit all-good and single-cancelling-challenge profiles, with their
    exact same-set bad events.
\end{description}
\end{leanledger}
The standalone import closure does not depend on HCP, protocol, consensus, or
runtime axiom bundles.
Both directions of the capacity-grid equality and its exact fractional-target
criterion are theorems, not axioms.  The next-shell lower bound is also
machine-checked.  The pair-MCA theorem and its \(2^{-128}\) family
instantiation are machine-checked as well.
The generic family corollaries retain their parameterized forms. The
realization module additionally constructs \(\mathbb F_{3^{n_u}}\), an
exact-order \(n_u\) root, and its injectively indexed subgroup for every
\(u\), and instantiates those corollaries without a realization hypothesis.
The standalone aggregate \lean{lean/ProximityCoding.lean} includes the
polynomial agreement modules. Its proof and ownership register is
\lean{papers/paper3/SEMANTIC_PARITY.md}. The 52 coding and 22 application
selections retain their target lists.
The \lean{Proximity/MCA/Syndrome} modules verify the projective incidence
criteria, the exterior-square test, the four-point maximum and the
unit-ideal certificate for the specified matching over \(\mathbb F_{17}\).
For that matching, \lean{Seventeen/Descent} checks the four distinct hits
over a field with a root of \(x^2=11\), and checks the Frobenius action on
the root. \lean{Seventeen/RadiusThree} checks the finite slope tables at
count three. The uniform spectrum in \cref{prop:seventeen-full-spectrum}
uses these certificates and the rank and counting arguments in the text.
The Gr\"obner procedure and marked point counts in
\cref{sec:mca-projective-incidence}, and the Zariski closure and tensor-price
arguments in \cref{sec:polynomial-agreement}, are given in the manuscript.
The support-intersection formula and stabilization in \cref{sec:list-descent}
are likewise mathematical arguments, using the checked ball equivalence.

\subsection{Auxiliary source and projective proofs}

The optional geometric source in
\cref{sec:projective-observation-boundary} uses the existing
\lean{Proximity/Concrete/GrassmannianPlucker/Reconstruction} module.
Its \lean{recoverSubspace_frame} and \lean{recoverHarmonicProjector_frame}
theorems recover the full subspace and its fixed-metric projector.
\lean{TCB/PluckerObservation.recover_eq_operativeHarmonicComponent}
adds the explicit operative-kernel identification.
The shared metric in \lean{Laplacian/Grassmannian/Metric} uses the Frobenius
distance of orthogonal projectors scaled by \(1/\sqrt{2}\).
Under \lean{Proximity/Concrete/GrassmannianPlucker}, the
\lean{ProjectorMetric} module proves that \lean{projectiveChordalDistance}
agrees with the original chordal distance on genuine frames and vanishes
exactly for equal shadows. The \lean{ProjectivePoint} module identifies
equality of \lean{frameProjectivePoint} values in Mathlib's projectivization
with equality of those shadows.
\lean{Proximity/Concrete/ProjectionChain/Core} composes the supplied
projective, cochain, plane and RS maps and their validity/near comparisons.
The exact-component obstruction in
\cref{prop:exact-component-obstruction} uses the existing
\lean{Sheaf/ObservationQuotient} proofs
\lean{HodgeEncryptedCochain.not_residual_norm_factors_exactComponent} and
\lean{HodgeEncryptedCochain.not_far_preserving_exactComponent}.
They retain the unrestricted Hodge carrier, a nonzero coexact direction,
exact-component factorization and zero target distance at the zero image.
The canonical systematic sheaf has \(\delta_1=0\), so this nonzero-coexact
premise is unavailable there.

These source-obstruction and projective results are canonical auxiliary proofs
outside the configured 52-target coding and 22-target native-modular
application selections. They do not enter the standalone coding import
closure; moving a statement into this manuscript does not add a selected
target. Their formal owners and the independent source metrics remain
distinct from Paper 2's codec, visibility and authentication contracts.
The projection-chain composition gives validity and near-side handoffs;
source-far preservation requires the separately stated calibration.

\pagebreak
\subsection{Status ledger}

\begingroup
\small
\begin{longtable}{@{}>{\raggedright\arraybackslash}p{0.42\linewidth}
  >{\raggedright\arraybackslash}p{0.20\linewidth}
  >{\raggedright\arraybackslash}p{0.28\linewidth}@{}}
\toprule
Obligation & Status & Required endpoint\\
\midrule
\endfirsthead
\toprule
Obligation & Status & Required endpoint\\
\midrule
\endhead
Count/relative ball correspondence & Lean-proved & Equality at \(r/n\),
  for \(n>0\)\\
Plane--pair equivalence & Lean-proved & Bijection and Hamming isometry\\
Plane--affine-line equivalence & Lean-proved & Formal inverse maps\\
Smooth large-field family & Lean-proved & Actual fields and subgroup domains\\
Capacity-grid list equality & Lean-proved & Universal \(L_m=\binom nk\)\\
Endpoint fractional criterion & Lean-proved & Target iff
  \(\binom nk\le\epsstar|\F|\)\\
Next-shell list lower bound & Lean-proved & Explicit \(|\F|^m\) injection\\
Complete list phase split & Lean-proved & Three regimes and exact real interval\\
One-codeword list budget & Lean-proved & Safe iff \(2r<n-k+1\)\\
Boundary-field list threshold & Lean-proved &
  Greatest safe count \(1536\); first unsafe count \(1537\)\\
Agreement-subset packing & Lean-proved & List bound at every count through capacity\\
List thresholds on Jo's length-\(64\) codes & Lean-proved & Four rates,
  certified prime fields and domains\\
Subfield field, domain, and arithmetic & Lean-proved & Order \(128\) and
  the budget comparisons, including list safety at count \(62\)\\
Subfield witnesses for both thresholds & Mathematical & Prefix-class and
  pole separation on the length-\(128\) code\\
Generic minimum-distance interpretation & Lean-proved & Attained minima,
  covering and good-challenge farness\\
Polynomial agreement certificates & Lean-proved & Agreement at most \(D\)
  and count bound \(n-D\)\\
Polynomial certificate degree boundary & Lean-proved & Exact normalized
  system, affine fibers, upward feasibility\\
Polynomial certificate completeness and sharpness & Lean-proved & Exact
  through half the minimum distance and sharp count \(\min(e,n-k+1-e)\)
  on nonzero domains\\
Varying-data closure and Euclidean prices & Mathematical & Elimination
  closure and tensor price bound\\
Pair-MCA certificate bound & Lean-proved &
  \(|\Bad_{\MCA}|\le\binom n{k+1}\)\\
MCA target-fraction inequality & Lean-proved &
  \(\binom n{k+1}\le\epsstar|\F|\) for the explicit family\\
Exact supplied-code MCA grid API & Lean-proved &
  Operational target iff \(2^{128}B_{\mathcal C}(r)\le|\F|\), with
  concrete next-line pinning\\
Exact RS small-radius certificates & Lean-proved &
  \(B_{\mathcal C}(0)=1\), and \(B_{\mathcal C}(1)\ge2\) for
  \(1\le k,\ k+2\le n\)\\
Smooth boundary-field instance & Lean-proved &
  Prime \(q=2^{128}+192513\equiv1\pmod {4096}\)\\
Boundary-field MCA threshold & Lean-proved &
  Greatest safe count \(0\); first unsafe count \(1\)\\
Explicit-family MCA greatest radius & Lean-proved &
  Full count radius \(n\) for the constructed family\\
Boundary-field real supremum & Lean-proved &
  \(1/4096\), not attained\\
Native scalar/pair metric comparison & Lean-proved &
  \(d_C\le d_P\le2d_C\)\\
Intrinsic modular application & Lean-proved & Actual integral coboundary
  validity, full original payload and \(d_C\le d_P=D_2\le2d_C\)\\
Sharp near scaling & Lean-proved & Factor two attained by \((u,0)\)\\
Canonical raw payload & Lean-proved & All three distances \(3/4\) at
  \((q,n,k)=(12289,4096,1024)\)\\
Monomial capacity far family & Lean-proved &
  Scalar, pair and flattened distances \((n-k)/n\)\\
Explicit challenge profiles & Lean-proved &
  All-good zero direction; diagonal bad set \(\{-1\}\)\\
Arbitrary fixed-mix obstruction & Lean-proved &
  Capacity-far pair with challenged word zero\\
Exact-component residual obstruction & Lean-proved, auxiliary &
  Nonfactorization and no positive far transfer under the stated
  unrestricted/nonzero-coexact hypotheses\\
Full projective reconstruction and metric & Lean-proved, auxiliary &
  Subspace/projector recovery, chordal separation and Mathlib projective-point
  agreement on genuine frames\\
\bottomrule
\end{longtable}
\endgroup

The counting, calibration, boundary-field and canonical application proofs
use only the standard Lean logical axioms
\lean{propext}, \lean{Classical.choice}, and \lean{Quot.sound}.
The arithmetic certificates are kernel-checked: the boundary primality proof
uses an explicit Lucas factor chain, and modular powers are checked by binary
modular-exponentiation proof terms. The field, roots and domain parameters
are unchanged. Earlier versions used \lean{native_decide} for some of these
certificates; their dated dependency reports are historical evidence rather
than the trust footprint of the present proofs. This arithmetic change does
not construct a commitment protocol or a stronger source-metric comparison.

\subsection{Verification criteria}

A completed claim requires:
\begin{enumerate}[label=(\roman*)]
  \item no placeholder assumptions for the main list/MCA result;
  \item an explicit statement of imported literature theorems;
  \item all field-size and rounding conditions checked;
  \item a clear separation between proved theorems, computational evidence, and
    conjectures;
  \item a reproducible proof artifact and manuscript cross-reference.
\end{enumerate}

\subsection{Remaining interfaces}

The nine named standalone coding statements now have complete Lean
endpoints, including the all-parameter smooth-family realization, complete
list phase split, real safe intervals and their excluded suprema, and the
generic minimum-distance interpretation. The generic source-to-pair calibration implication
\[
  d_{\mathrm{source}}(x,\mathrm{valid})\ge\delta
  \Longrightarrow
  \Delta(\pi(x),\mathcal C^{\equiv2})\ge\delta
\]
is Lean-proved for an abstract distance model from its named pointwise
non-collapse field. The distance-calibration module identifies the target
model with the generic attained pair-code minimum. Pullback distance
makes that field reflexive; an application choosing a different pre-existing
source metric supplies one explicit comparison.
The one-codeword and immediately sub-capacity budgets have exact
thresholds. For prescribed domains with intermediate budgets below the
packing cut, the general optimum remains open.

The concrete application hierarchy contains the extraction, plane/pair
Hamming isometry, distance equality and pullback calibration of Section~7.
These constructions lie outside the standalone \lean{ProximityCoding}
import closure. Section~7 instantiates them for the canonical modular
profile. The restricted native Hodge comparison in
\cref{prop:restricted-hodge-rs} uses its own structured samples and a
sitewise error bound. Other HCP or frame scores require their own
source comparison.

The explicit source realization lives in the \lean{Sheaf/SystematicCode}
modules, including \lean{Modular}. Their actual integral cellular coboundary
is proved compatible with RS reduction. The storage modules
\lean{CoordinateCodec}, \lean{NoWrapCalibration} and \lean{ThreeLevelMetrics}
retain the original payload and separate the restricted Euclidean comparison
from the global native-metric comparison. The canonical specialization and
public-observation joins are in \lean{CanonicalSystematic} and
\lean{CanonicalThreeLevel}; \lean{CanonicalCalibration} composes the existing
interfaces with their declared near/far scales. These modules live under
\lean{Protocol/Storage}. Their file location does not change manuscript
ownership: Paper 2 specifies the shared construction and source/observation
contracts \cite{vanDerVlugt2026protocol}; this paper proves the coding and
quantitative metric comparisons for that exact construction. The two Paper 3
verification configurations separate core and bridge statements for source-size
purposes, without introducing another paper or protocol result.

These explicit profiles do not identify arbitrary old HCP presentations or
their frame scores with native modular edit distance. A stronger upstream
geometric comparison, or a different challenge/observation profile, needs its
own theorem. Such obligations remain separate from the completed generic
coding results and do not require Paper~2 consensus or Rust.
